\chapter{第9章：Galois 表示 习题}
\difficultykey

\section{基础题 \easy}

\begin{problem}{Krull拓扑的基本邻域}
\easy
设
\[
G_\Q=\Gal(\bar\Q/\Q).
\]
说明 Krull 拓扑的一组邻域基可以取为
\[
\Gal(\bar\Q/L),
\]
其中 $L/\Q$ 跑遍有限 Galois 扩张；并说明这些子群既开又闭。
\end{problem}
\begin{solution}
Krull 拓扑按定义正是使所有限制映射
\[
G_\Q\longrightarrow \Gal(L/\Q)
\]
连续的最弱拓扑，其中 $L/\Q$ 取有限 Galois 扩张。于是单位元的基本开邻域正是这些映射的核，即
\[
U_L=\Gal(\bar\Q/L).
\]
任一点 $\sigma\in G_\Q$ 的基本邻域便是陪集 $\sigma U_L$。

又因为
\[
G_\Q/U_L\cong \Gal(L/\Q)
\]
是有限离散群，所以商映射连续，$U_L$ 是开子群；有限指数开子群的补是有限个陪集之并，因此也是闭集。故 $U_L$ 既开又闭。
\end{solution}

\begin{problem}{连续性的等价刻画}
\easy
设
\[
\rho:G_\Q\to \GL_n(\Z_\ell)
\]
是一个群同态。证明下列条件等价：
\begin{itemize}
\item[(i)] $\rho$ 连续；
\item[(ii)] 对每个 $m\ge1$，复合同态
\[
G_\Q\xrightarrow{\rho}\GL_n(\Z_\ell)\to \GL_n(\Z/\ell^m\Z)
\]
都有开核；
\item[(iii)] 对每个 $m\ge1$，$\rho^{-1}(1+\ell^mM_n(\Z_\ell))$ 都是开子群。
\end{itemize}
\end{problem}
\begin{solution}
在 $\GL_n(\Z_\ell)$ 中，主同余子群
\[
K_m=1+\ell^mM_n(\Z_\ell)
\]
构成单位元的一组邻域基，所以 $\rho$ 连续当且仅当每个 $\rho^{-1}(K_m)$ 开，这就是 (i)$\Leftrightarrow$(iii)。

而 $K_m$ 正是约化映射
\[
\GL_n(\Z_\ell)\to \GL_n(\Z/\ell^m\Z)
\]
的核，因此复合同态的核恰等于 $\rho^{-1}(K_m)$，故 (ii) 与 (iii) 等价。
\end{solution}

\begin{problem}{分圆特征的定义}
\easy
设 $\ell$ 为素数。说明 Galois 群对所有 $\ell^n$ 次单位根的作用给出同态
\[
\chi_\ell:G_\Q\to \Z_\ell^\times,
\]
并证明对每个 $n$，其模 $\ell^n$ 约化等于
\[
G_\Q\to \Aut(\mu_{\ell^n})\cong (\Z/\ell^n\Z)^\times.
\]
\end{problem}
\begin{solution}
对任意 $\sigma\in G_\Q$ 与 $\zeta\in\mu_{\ell^n}$，有
\[
\sigma(\zeta)=\zeta^{a_n(\sigma)}
\]
其中 $a_n(\sigma)\in(\Z/\ell^n\Z)^\times$ 唯一确定。这给出同态
\[
\chi_{\ell,n}:G_\Q\to(\Z/\ell^n\Z)^\times.
\]
当 $m<n$ 时，$\mu_{\ell^m}\subset\mu_{\ell^n}$，故这些同态彼此相容，于是取逆极限得到
\[
\chi_\ell=\varprojlim_n\chi_{\ell,n}:G_\Q\to\varprojlim_n(\Z/\ell^n\Z)^\times=\Z_\ell^\times.
\]
按构造，$\chi_\ell$ 模 $\ell^n$ 的约化正是 $\chi_{\ell,n}$。
\end{solution}

\begin{problem}{二挠点的直接计算}
\easy
对椭圆曲线
\[
E:y^2=x^3-x
\]
写出 $E[2]$ 的全部点，并说明 $G_\Q$ 在 $E[2]$ 上的作用为何是平凡的。
\end{problem}
\begin{solution}
把右边分解为
\[
x^3-x=x(x-1)(x+1).
\]
非零二挠点正是 $y=0$ 时的三根点，因此
\[
E[2]=\{\mathcal O,(0,0),(1,0),(-1,0)\}.
\]
这些点的坐标都在 $\Q$ 中，所以对任何 $\sigma\in G_\Q$ 都有
\[
\sigma(0,0)=(0,0),\quad \sigma(1,0)=(1,0),\quad \sigma(-1,0)=(-1,0).
\]
故 $G_\Q$ 在 $E[2]$ 上逐点固定，因而该作用是平凡的。
\end{solution}

\begin{problem}{Tate模的相容序列}
\easy
说明 Tate 模
\[
T_\ell(E)=\varprojlim_n E[\ell^n]
\]
的元素就是满足
\[
\ell P_{n+1}=P_n
\]
的相容序列 $(P_n)_{n\ge1}$；并说明在特征零时它是自由秩 $2$ 的 $\Z_\ell$-模。
\end{problem}
\begin{solution}
逆极限的定义就是
\[
T_\ell(E)=\left\{(P_n)\in\prod_{n\ge1}E[\ell^n]:\ell P_{n+1}=P_n\text{ 对所有 }n\right\}.
\]
因此它的元素正是不断向上提升的挠点链。

在特征零时，$E[\ell^n]$ 有 $\ell^{2n}$ 个点，有限 Abel 群分类给出
\[
E[\ell^n]\cong (\Z/\ell^n\Z)^2.
\]
对 $n$ 取逆极限便得
\[
T_\ell(E)\cong \varprojlim_n(\Z/\ell^n\Z)^2\cong \Z_\ell^2.
\]
所以它是自由秩 $2$ 的 $\Z_\ell$-模。
\end{solution}

\begin{problem}{导子十一曲线的前两个系数}
\easy
设
\[
E_{11}:y^2+y=x^3-x^2-10x-20.
\]
直接数点计算 $a_2(E_{11})$ 与 $a_3(E_{11})$。
\end{problem}
\begin{solution}
模 $2$ 时方程化为
\[
y^2+y=x^3+x^2.
\]
代入 $x=0,1$ 都得到右边为 $0$，而 $y=0,1$ 都满足 $y^2+y=0$。故
\[
\#E_{11}(\F_2)=5,
\qquad
 a_2=2+1-5=-2.
\]

模 $3$ 时方程化为
\[
y^2+y=x^3-x^2+2x+1.
\]
当 $x=0$ 时右边为 $1$，无解；当 $x=1,2$ 时右边都为 $0$，各给出两个解，再加无穷远点，所以
\[
\#E_{11}(\F_3)=5,
\qquad
 a_3=3+1-5=-1.
\]
\end{solution}

\begin{problem}{Frobenius特征多项式与平方域点数}
\easy
设椭圆曲线 $E/\Q$ 在素数 $p\neq \ell$ 处好约化，且 $\Frob_p$ 在 $T_\ell(E)$ 上的特征多项式为
\[
X^2-a_pX+p.
\]
设其根为 $\alpha,\beta$。证明
\[
\#E(\F_{p^2})=p^2+1-(\alpha^2+\beta^2).
\]
并用 $E_{11}$ 在 $p=2$ 时验证结果。
\end{problem}
\begin{solution}
由 Weil 猜想在椭圆曲线情形的标准公式，
\[
\#E(\F_{p^r})=p^r+1-(\alpha^r+\beta^r).
\]
取 $r=2$ 即得
\[
\#E(\F_{p^2})=p^2+1-(\alpha^2+\beta^2).
\]

对 $E_{11}$，上题算得 $a_2=-2$，故特征多项式为
\[
X^2+2X+2.
\]
于是
\[
\alpha+\beta=-2,
\qquad
\alpha\beta=2,
\]
从而
\[
\alpha^2+\beta^2=(\alpha+\beta)^2-2\alpha\beta=4-4=0.
\]
因此
\[
\#E_{11}(\F_4)=4+1-0=5.
\]
\end{solution}

\begin{problem}{行列式与分圆特征}
\easy
证明对椭圆曲线的 $\ell$-进表示有
\[
\det\rho_{E,\ell}=\chi_\ell.
\]
并说明在好素数 $p\neq\ell$ 处这给出
\[
\det\rho_{E,\ell}(\Frob_p)=p.
\]
\end{problem}
\begin{solution}
Weil 配对给出一个 Galois 相容的双线性配对
\[
e_\ell:T_\ell(E)\times T_\ell(E)\to\Z_\ell(1).
\]
若取 $T_\ell(E)$ 的一组基，Galois 元 $\sigma$ 在左边的作用会使配对乘上矩阵行列式，而在右边的作用正是乘以分圆特征 $\chi_\ell(\sigma)$。因此
\[
\det\rho_{E,\ell}(\sigma)=\chi_\ell(\sigma)
\]
对所有 $\sigma\in G_\Q$ 都成立。

当 $p\neq\ell$ 且表示在 $p$ 处未分歧时，算术 Frobenius 在 $\ell$ 次单位根上的作用是 $\zeta\mapsto\zeta^p$，故
\[
\chi_\ell(\Frob_p)=p.
\]
于是
\[
\det\rho_{E,\ell}(\Frob_p)=p.
\]
\end{solution}

\begin{problem}{模表示的约化}
\easy
设
\[
\rho:G_\Q\to\GL_2(\Z_\ell)
\]
连续。说明把矩阵项逐项模 $\ell$ 约化后得到表示
\[
\bar\rho:G_\Q\to\GL_2(\F_\ell),
\]
并解释为什么它只依赖于 $\rho$ 的同构类而不依赖于所选的 $\Z_\ell$-基。
\end{problem}
\begin{solution}
由约化映射
\[
\GL_2(\Z_\ell)\to\GL_2(\F_\ell)
\]
复合可得
\[
\bar\rho:G_\Q\to\GL_2(\F_\ell).
\]
因为矩阵乘法与约化相容，所以这仍是群同态。

若换一组 $\Z_\ell$-基，则 $\rho$ 会被某个 $A\in\GL_2(\Z_\ell)$ 共轭为
\[
\rho'(\sigma)=A\rho(\sigma)A^{-1}.
\]
模 $\ell$ 约化后得到
\[
\bar\rho'(\sigma)=\bar A\,\bar\rho(\sigma)\,\bar A^{-1}.
\]
因此新旧约化表示彼此同构，所以 $\bar\rho$ 只依赖于原表示的同构类。
\end{solution}

\section{进阶题 \medium}

\begin{problem}{有理挠点与可约性}
\medium
设椭圆曲线 $E/\Q$ 具有一个有理 $\ell$-挠点 $P\in E(\Q)$。证明残余表示
\[
\bar\rho_{E,\ell}:G_\Q\to\GL_2(\F_\ell)
\]
是可约的。
\end{problem}
\begin{solution}
因为 $P$ 是有理点，对任意 $\sigma\in G_\Q$ 都有 $\sigma(P)=P$。于是 $E[\ell]$ 中由 $P$ 张成的一维子空间
\[
\F_\ell P\subset E[\ell]
\]
在 Galois 作用下稳定。

二维表示存在非零真不变子空间正是“可约”的定义，所以 $\bar\rho_{E,\ell}$ 可约。
\end{solution}

\begin{problem}{二进残余表示的判别}
\medium
证明对椭圆曲线 $E/\Q$，表示
\[
\bar\rho_{E,2}:G_\Q\to\GL_2(\F_2)
\]
可约当且仅当 $E$ 有非平凡有理二挠点。
\end{problem}
\begin{solution}
若 $E$ 有非平凡有理二挠点 $P$，上题立刻给出 $\bar\rho_{E,2}$ 可约。

反过来，若 $\bar\rho_{E,2}$ 可约，则二维 $\F_2$-空间 $E[2]$ 中存在一维 Galois 不变子空间。由于 $\F_2$ 只有一个非零元，该一维子空间恰对应某个非零点 $P\in E[2]$，并且对所有 $\sigma$ 都满足 $\sigma(P)=P$。因此 $P\in E(\Q)$，即存在非平凡有理二挠点。
\end{solution}

\begin{problem}{可约残余表示的迹同余}
\medium
设某个二维模 $\ell$ 表示满足半单化分解
\[
\bar\rho^{\mathrm{ss}}\cong 1\oplus\bar\chi_\ell.
\]
证明对所有未分歧于该表示的素数 $p\neq\ell$，都有
\[
\Tr\,\bar\rho(\Frob_p)\equiv 1+p\pmod \ell.
\]
\end{problem}
\begin{solution}
半单化分解给出
\[
\bar\rho^{\mathrm{ss}}(\Frob_p)\sim
\begin{pmatrix}
1&0\\0&\bar\chi_\ell(\Frob_p)
\end{pmatrix}.
\]
因此迹等于两条对角线之和，即
\[
\Tr\,\bar\rho(\Frob_p)=1+\bar\chi_\ell(\Frob_p).
\]
而对 $p\neq\ell$ 的 Frobenius，分圆特征模 $\ell$ 的值是 $p\bmod\ell$，所以
\[
\Tr\,\bar\rho(\Frob_p)\equiv 1+p\pmod\ell.
\]
\end{solution}

\begin{problem}{五挠点给出的同余}
\medium
已知导子 $11$ 的椭圆曲线 $E_{11}$ 的有理挠子群同构于 $\Z/5\Z$。证明对每个好素数 $p\neq5,11$，都有
\[
a_p(E_{11})\equiv 1+p\pmod5,
\]
并检验 $p=2,3$ 的情形。
\end{problem}
\begin{solution}
因为 $E_{11}$ 有有理 $5$-挠点，所以上题表明
\[
\bar\rho_{E_{11},5}
\]
可约。其行列式是 $\bar\chi_5$，故半单化必为
\[
1\oplus\bar\chi_5.
\]
于是对所有好素数 $p\neq5,11$，有
\[
a_p(E_{11})\equiv \Tr\,\bar\rho_{E_{11},5}(\Frob_p)\equiv 1+p\pmod5.
\]

上面已算出 $a_2=-2$ 与 $a_3=-1$。模 $5$ 化后
\[
-2\equiv3\equiv1+2,
\qquad
-1\equiv4\equiv1+3.
\]
检验无误。
\end{solution}

\begin{problem}{Deligne表示与椭圆曲线表示}
\medium
设 $f=\sum_{n\ge1}a_nq^n$ 是权 $2$、级别 $N$ 的归一化新形式，对应椭圆曲线 $E_f/\Q$。说明为何对每个好素数 $p\nmid N\ell$，表示 $\rho_f$ 与 $\rho_{E_f,\ell}$ 在 $\Frob_p$ 上具有同一特征多项式，并解释这如何推出
\[
\rho_f\cong\rho_{E_f,\ell}.
\]
\end{problem}
\begin{solution}
Deligne 的表示满足
\[
\det(1-X\rho_f(\Frob_p))=1-a_pX+pX^2,
\]
而椭圆曲线的 Tate 模表示满足
\[
\det(1-X\rho_{E_f,\ell}(\Frob_p))=1-a_p(E_f)X+pX^2.
\]
由 Eichler--Shimura，$a_p(E_f)=a_p(f)$，所以两者在所有好素数处的特征多项式一致。

两表示都只在有限多素数处分歧，又都是二维半单表示。由 Chebotarev 密度定理，Frobenius 共轭类在未分歧处稠密，因而这些特征多项式就决定了半单表示。故
\[
\rho_f\cong\rho_{E_f,\ell}.
\]
\end{solution}

\begin{problem}{Chebotarev的基本应用}
\medium
设
\[
\rho_1,\rho_2:G_\Q\to\GL_2(\Q_\ell)
\]
是两个半单连续表示，只在有限多素数处分歧。若对所有充分大的素数 $p$ 都有
\[
\Tr\,\rho_1(\Frob_p)=\Tr\,\rho_2(\Frob_p),
\qquad
\det\rho_1(\Frob_p)=\det\rho_2(\Frob_p),
\]
证明 $\rho_1\cong\rho_2$。
\end{problem}
\begin{solution}
在二维情形，迹与行列式决定特征多项式，所以假设说明两表示在所有充分大的未分歧素数处有相同的 Frobenius 特征多项式。

Chebotarev 密度定理告诉我们，这些 Frobenius 共轭类在去掉有限分歧集后是稠密的。对半单 $\ell$-进表示，Brauer--Nesbitt 型结论说明：若两个表示在稠密集合上有相同特征多项式，则二者同构。因此
\[
\rho_1\cong\rho_2.
\]
\end{solution}

\begin{problem}{变基与共轭不变量}
\medium
说明若把 $T_\ell(E)$ 的基从 $\mathcal B$ 换为 $\mathcal B'$，则表示矩阵只会被同一个矩阵共轭。因此
\[
\Tr\rho_{E,\ell}(\sigma),\qquad \det\rho_{E,\ell}(\sigma)
\]
都与基无关。
\end{problem}
\begin{solution}
设从基 $\mathcal B$ 到 $\mathcal B'$ 的换基矩阵为 $A\in\GL_2(\Z_\ell)$。若在基 $\mathcal B$ 下 $\sigma$ 的矩阵为 $M(\sigma)$，则在线性代数中同一线性变换在新基下的矩阵为
\[
M'(\sigma)=AM(\sigma)A^{-1}.
\]
所以换基只会导致整体共轭。

而迹与行列式都是共轭不变量，即
\[
\Tr(AMA^{-1})=\Tr(M),
\qquad
\det(AMA^{-1})=\det(M).
\]
故它们与所选基无关，只由表示本身决定。
\end{solution}

\begin{problem}{水平下降定理的陈述}
\medium
在半稳定椭圆曲线的情形下，陈述 Ribet 水平下降定理的常用版本，并解释其输入与输出分别是什么。
\end{problem}
\begin{solution}
常用的陈述是：设 $E/\Q$ 为半稳定椭圆曲线，导子为 $N$，$\ell$ 为素数。若残余表示
\[
\bar\rho_{E,\ell}:G_\Q\to\GL_2(\F_\ell)
\]
不可约，并且它来自某个权 $2$、级别 $N$ 的新形式，那么 $\bar\rho_{E,\ell}$ 也来自较低级别的新形式；在 Frey 曲线应用中，所有与乘法约化对应的坏素数都可以从级别中去掉，最终降到很小的级别。

输入是：半稳定性、模性、以及残余表示不可约。输出是：同一个模表示可以在更低的级别上出现，即“水平下降”。
\end{solution}

\begin{problem}{局部因子的一致性}
\medium
设椭圆曲线 $E/\Q$ 与归一化新形式 $f$ 满足对所有好素数 $p\nmid N\ell$ 都有
\[
\det(1-X\rho_{E,\ell}(\Frob_p))=
\det(1-X\rho_f(\Frob_p)).
\]
说明它们在这些素数处的 Euler 因子完全相同，因此对应的 $L$-函数只可能在有限多个坏素数处不同。
\end{problem}
\begin{solution}
椭圆曲线在好素数 $p$ 处的 Euler 因子为
\[
L_p(E,s)^{-1}=1-a_p(E)p^{-s}+p^{1-2s},
\]
而模形式表示的 Euler 因子为
\[
L_p(f,s)^{-1}=1-a_p(f)p^{-s}+p^{1-2s}.
\]
给定特征多项式一致，就说明
\[
a_p(E)=a_p(f)
\]
且常数项同为 $p$，所以两者局部因子相同。

因此在所有共同的好素数处，Euler 乘积逐项一致。只有在有限多个坏素数处，局部因子可能需要单独比较。
\end{solution}

\section{挑战题 \hard}

\begin{problem}{稠密Frobenius与表示唯一性}
\hard
设
\[
\rho_1,\rho_2:G_\Q\to\GL_2(\Q_\ell)
\]
是两个半单连续表示，只在有限多素数处分歧。若对所有与分歧集不相交的素数 $p$，都有相同的 Frobenius 特征多项式，说明为什么这已经足以唯一确定半单表示。
\end{problem}
\begin{solution}
每个未分歧素数 $p$ 给出一个 Frobenius 共轭类 $\Frob_p$。Chebotarev 密度定理说明，这些共轭类在去掉有限分歧集后形成稠密子集。

若两表示在全部这些共轭类上都有相同的特征多项式，则它们的特征函数在一个稠密集合上一致。对半单 $\ell$-进表示，Brauer--Nesbitt 原理告诉我们，角色相同就意味着表示同构。因此 Frobenius 特征多项式在“几乎所有 $p$”上的数据已经决定了半单表示。
\end{solution}

\begin{problem}{所有迹都为零的不可能性}
\hard
说明不存在椭圆曲线 $E/\Q$ 满足对所有好素数 $p$ 都有
\[
a_p(E)=0.
\]
提示：把这些迹看成二维 Galois 表示的迹。
\end{problem}
\begin{solution}
若对所有好素数 $p\neq\ell$ 都有 $a_p(E)=0$，则 $\rho_{E,\ell}(\Frob_p)$ 的迹在一个稠密的 Frobenius 集合上恒为零。迹作为矩阵项的和是连续函数，所以它在该稠密集上的值决定了整个表示的迹函数；于是应有
\[
\Tr\rho_{E,\ell}(\sigma)=0
\]
对所有 $\sigma\in G_\Q$ 都成立。

但对恒等元 $1\in G_\Q$，矩阵是单位矩阵，故
\[
\Tr\rho_{E,\ell}(1)=2,
\]
矛盾。因此不可能所有好素数处都满足 $a_p(E)=0$。
\end{solution}

\begin{problem}{水平下降的数学含义}
\hard
设半稳定模椭圆曲线 $E/\Q$ 的导子为 $N$，某个素数 $q\mid N$ 对应乘法约化。解释“把级别中的 $q$ 去掉”在 Galois 表示上意味着什么。
\end{problem}
\begin{solution}
导子中的素因子 $q$ 记录了表示在 $q$ 处的 ramification。半稳定且在 $q$ 处乘法约化时，$\ell$-进表示在 $q$ 处并非完全平凡，但其导子指数只有 $1$。

若 Ribet 水平下降说明同一个残余表示也来自级别 $N/q$ 的新形式，这意味着模 $\ell$ 以后，原先在 $q$ 处导致导子出现的局部 ramification 已经“消失”了；换言之，$\bar\rho_{E,\ell}$ 可以由一个在 $q$ 处更简单的模形式实现。数学上，水平下降就是把残余表示的导子变小。
\end{solution}

\begin{problem}{有限像表示的确定}
\hard
设
\[
\rho:G_\Q\to\GL_n(\C)
\]
的像有限。证明：若已知对几乎所有素数 $p$ 的共轭类 $\rho(\Frob_p)$，那么 $\rho$ 的角色就被唯一确定。
\end{problem}
\begin{solution}
像有限意味着 $\ker\rho$ 是开正规子群，所以对应一个有限 Galois 扩张 $L/\Q$，使 $\rho$ 通过有限群
\[
\Gal(L/\Q)
\]
分解。

Chebotarev 定理对有限扩张 $L/\Q$ 说：$\Gal(L/\Q)$ 的每个共轭类都会以正密度出现为某个未分歧素数的 Frobenius 共轭类。因此只要知道几乎所有 $p$ 的 $\rho(\Frob_p)$，就等于知道有限像群中每个共轭类上的角色值。角色决定复半单表示，所以 $\rho$ 被唯一确定。
\end{solution}

\begin{problem}{Serre猜想的陈述}
\hard
陈述 Serre 关于二维奇不可约模 $\ell$ Galois 表示的猜想，并说明它与本章内容的关系。
\end{problem}
\begin{solution}
Serre 猜想断言：任意连续、奇、不可约的表示
\[
\bar\rho:G_\Q\to\GL_2(\overline{\F}_\ell)
\]
都来自某个模形式；更精细地说，还能预测最小权与最小级别。

它与本章的关系是：本章从椭圆曲线或模形式出发构造 Galois 表示，而 Serre 猜想给出反方向的纲领——从模表示反推模形式。Ribet 水平下降正是这一纲领中的关键环节之一。
\end{solution}

\begin{problem}{同一表示与Fourier系数}
\hard
设两个归一化新形式 $f,g$ 的 $\ell$-进表示同构：
\[
\rho_f\cong\rho_g.
\]
证明对所有不整除 $N_fN_g\ell$ 的素数 $p$，都有
\[
a_p(f)=a_p(g).
\]
并解释为何这强烈暗示 $f=g$。
\end{problem}
\begin{solution}
表示同构意味着对每个这样的好素数 $p$，Frobenius 的特征多项式相同。Deligne 表示的特征多项式为
\[
X^2-a_p(f)X+p,
\qquad
X^2-a_p(g)X+p,
\]
故立刻得到
\[
a_p(f)=a_p(g).
\]

对归一化 Hecke 本征形式而言，几乎所有素数处的 Hecke 特征值就决定整条本征形式；这就是 multiplicity one 现象。因此同一 Galois 表示通常只能对应同一个归一化新形式。
\end{solution}

\begin{problem}{残余分解与一般同余}
\hard
设某二维模 $\ell$ 表示的半单化分解为
\[
\bar\rho^{\mathrm{ss}}\cong \varphi\oplus \bar\chi_\ell\varphi^{-1},
\]
其中 $\varphi$ 是一维表示。证明对每个未分歧素数 $p\neq\ell$，都有
\[
\Tr\,\bar\rho(\Frob_p)
\equiv
\varphi(\Frob_p)+p\,\varphi(\Frob_p)^{-1}
\pmod\ell.
\]
\end{problem}
\begin{solution}
由半单化分解，$\Frob_p$ 在该表示上的特征值就是
\[
\varphi(\Frob_p),
\qquad
\bar\chi_\ell(\Frob_p)\varphi(\Frob_p)^{-1}.
\]
因此迹等于二者之和：
\[
\Tr\,\bar\rho(\Frob_p)=
\varphi(\Frob_p)+\bar\chi_\ell(\Frob_p)\varphi(\Frob_p)^{-1}.
\]
而 $\bar\chi_\ell(\Frob_p)\equiv p\pmod\ell$，所以得到所需同余
\[
\Tr\,\bar\rho(\Frob_p)
\equiv
\varphi(\Frob_p)+p\,\varphi(\Frob_p)^{-1}
\pmod\ell.
\]
\end{solution}

\begin{problem}{半稳定条件与水平下降假设}
\hard
解释为什么在半稳定椭圆曲线的应用中，Ribet 定理的局部条件特别容易核查。
\end{problem}
\begin{solution}
半稳定意味着每个坏素数处只有乘法约化，没有加性约化。因此局部导子指数最简单：好约化时为 $0$，乘法约化时为 $1$。这正是水平下降最容易处理的局部形状。

换言之，半稳定把复杂的局部可能性大量排除掉，使残余表示在坏素数处的行为足够简单，从而更容易证明这些素数在模 $\ell$ 情形下可以从级别中去掉。这就是 Frey 曲线和 Wiles 论证都偏爱半稳定对象的原因。
\end{solution}

\begin{problem}{级别二的最终矛盾}
\hard
设某个半稳定椭圆曲线 $E/\Q$ 已知是模的，并且对某个素数 $\ell$，残余表示 $\bar\rho_{E,\ell}$ 不可约。若 Ribet 定理把它的级别降到 $2$，说明为什么会立刻得到矛盾。
\end{problem}
\begin{solution}
若水平下降成功，则 $\bar\rho_{E,\ell}$ 必来自某个权 $2$、级别 $2$ 的新形式。于是必须存在非零空间
\[
S_2(\Gamma_0(2)).
\]
但前面章节已知
\[
\dim S_2(\Gamma_0(2))=0.
\]
因此根本不存在这样的新形式，更不可能有对应的残余表示。

所以“降到级别 $2$”本身就是矛盾。这正是 Fermat 大定理证明中最关键的终点之一。
\end{solution}
